\section{Introduction}
\label{sec:intro}

Strong gravitational lenses are rare. In Euclid's Q1 release, the expert-graded catalogue
holds a few hundred confident systems across 63 square degrees, about one for every 2200
objects searched~\citep{lines2025simstosky}. The objects a search has to find are scarce,
and so are the labelled examples a supervised classifier would need. The usual way around
this is to simulate lenses, train on the simulations and apply the model to the sky. That
turns a labelling problem into a distribution-shift problem: simulated arcs, noise and
deflectors differ from real ones, and in a ground-based search for lensing by late-type
galaxies a simulated test set overestimated the purity of the selected sample by at least a
factor of two~\citep{unions2025}.

We study that shift on Euclid Q1. We train lens finders on simulated arcs painted onto real
Euclid galaxies and adapt them toward real images with unsupervised domain adaptation. No
real-lens label is ever a training target. Labels do enter in other ways. We build each
fold's unlabelled target pool from catalogued lenses and random galaxies, and we balance
every target batch 1:1 using the catalogue labels (Section~\ref{sec:labels}). Because that
balancing is the design's main weakness, we test it in a second pre-registered experiment, with
a target diluted to one lens per ten galaxies whose batches are drawn without the labels
(Section~\ref{sec:results-dilute}). The Euclid Q1
sim-to-real benchmark~\citep{lines2025simstosky} reports a simulation-only lens finder and
invites adaptation follow-up work on the same data; this paper is such a follow-up.

We measure performance where a search operates: the true-positive rate at a false-positive
rate of $10^{-3}$. We run four published adaptation objectives on identical folds against a
control that fine-tunes the same encoder with its adaptation weight set to zero. The
evaluation positives are the 489 graded real lenses, five folds of 98 (one of 97), and that
sample size bounds every comparison we make.

Our contributions are these.
\begin{itemize}
\item \textbf{A controlled measurement of what a target-domain training signal buys.} A
      control that differs from the DANN arm only in its adaptation weight separates having
      a target signal from having none: on the class-balanced target of our main design, DANN
      raises TPR at FPR${}=10^{-3}$ by $\dCtDannLam{}$ over it, and the other three objectives,
      for which it is not a matched control, by $\dCtBnmLamAbs{}$ to $\dCtCoralLamPAbs{}$ (Section~\ref{sec:results-ablation}).
\item \textbf{A pre-registered test of the balanced target.} With the target diluted to one lens
      per ten galaxies and its batches drawn without the labels, DANN and BNM keep a detectable
      gain over the control; for CORAL and FreeMatch we cannot detect a gain.
      Dilution also changes how many lenses each batch holds and how often each is seen, so
      this does not isolate the class proportion (Section~\ref{sec:results-dilute}).
\item \textbf{How far the methods can be ranked.} At each arm's declared read-off, every paired
      interval among CORAL, DANN and FreeMatch includes zero. Averaged over \rwFiveKWord{}
      checkpoints, FreeMatch leads under each of two pre-declared score combinations, and the
      lead excludes zero under one of them; we treat it as suggestive. Settling the ranking
      needs more graded lenses rather than another read-off convention
      (Section~\ref{sec:resolution}).
\item \textbf{What the models find and miss.} In every model, recall is about two to three
      times higher for grade-A than for grade-B lenses. The highest-scoring negatives are mostly
      bright galaxies with halos, rings or spiral arms, and $\rvTopNegC{}$ of the
      $\rvTopNegN{}$ negatives that reach the top 30 of any of the seven readouts on any fold are grade-C
      candidates in the catalogue (Section~\ref{sec:findmiss}). Catalogued lenses hidden among
      the original test negatives cost the stronger models the most
      (Section~\ref{sec:testpool} and Appendix~\ref{app:prereg_pool}).
\end{itemize}

Sections~\ref{sec:data}--\ref{sec:eval} describe the data, the methods and the evaluation;
Section~\ref{sec:results} gives the results, with the test of the balanced target in
Section~\ref{sec:results-dilute}, Section~\ref{sec:negative} what did not work,
and Section~\ref{sec:limitations} the limitations, beginning with the target and its labels.
